\documentclass[12pt]{article}
\usepackage[margin=1in]{geometry}
\usepackage{amsmath, amssymb}
\usepackage{enumitem, array}

\usepackage{boxedminipage}
\usepackage{fancyhdr}
\usepackage{ragged2e, comment}


\newenvironment{problem}[1][]{%
  \bigskip% put space before problem statement box %
  \noindent\begin{boxedminipage}{\columnwidth}%
  \def\@tempa{#1}%
  \ifx\@tempa\empty\else%
    \textbf{#1}\hspace{0.5em}%
  \fi%
}{%
  \end{boxedminipage}%
}


\begin{document}

\begin{center}
    \Large \textbf{MATH 2790 Test 1 Fall 2025} \\[1em]
\end{center}

\vspace{.25cm}

\noindent Name: \underline{\hspace{8cm}} \hspace{.25cm} Auburn ID: \underline{\hspace{4cm}}



\section*{Instructions}
\begin{enumerate}[label=\arabic*.]
\item Print your full name and your banner ID on your exam. Then finish reading these instructions, and sign the bottom of this page.
\item All electronic equipment (cell phones, Apple Watch, AirPod, etc) are \textbf{NOT ALLOWED} during the exam time.
\item You are allowed to use up to 2 calculators from the following list: BA-35, BA II Plus, BA II Plus Professional, TI-30Xa, TI-30X II (IIS solar or IIB battery), TI-30XS MultiView (or XB battery).
\item No talking is allowed during the exam.
\item Unless otherwise indicated, \textbf{SHOW ALL YOUR WORK}. If no work is shown, no partial credit can be awarded. Your work and answer need to be accurate and relevant to receive points.
\item You will be given exactly 50 minutes for this exam.
\item Any student not following the above instructions nor behaving according to the above instructions during the exam may have their exam confiscated and points deducted.
\end{enumerate}

I have read and fully understand all of the above instructions. \\[2em]
Signature: \hrulefill


\vspace{.75cm}
Do not write in the table below

% scoring table
\begin{center}
\begin{tabular}{|*{3}{>{\centering\arraybackslash}p{.1\linewidth}|}}
    \hline
    \textbf{Problem} & \textbf{Points} & \textbf{Score} \\
    \hline
    1 & 10 &  \\
    \hline
    2 & 10 &  \\
    \hline
    3 & 15 &  \\
    \hline
    4 & 12 &  \\
    \hline
    5 & 16 &  \\
    \hline
    6 & 12 &  \\
    \hline
    7 & 12 &  \\
    \hline
    8 & 12 &  \\
    \hline \hline
    \textbf{Total} & 99+1 &  \\
    \hline
\end{tabular}
\end{center}

\newpage

%%%%% PROBLEM 1 %%%%%
\newpage
\begin{problem}[1. (10 points)]
Find the accumulated value of $\$5,000$ at the end of $8$ years if the nominal rate of discount was $4.5\%$ convertible quarterly for the first $1$ years, the effective monthly interest rate was $.025\%$ for the next $4$ years and the nominal interest rate was $2.5\%$ convertible continuously for the last $3$ years. Round to $4$ decimal places.
\end{problem}

\newpage
%%%%% PROBLEM 2 %%%%%
\newpage
\begin{problem}[2. (10 points)]
Owen took out a loan governed by simple interest at a rate of $6.5\%$. Find the discount rate $d_{[4,6]}$. Round to $4$ decimal places.
\end{problem}

\newpage

%%%%% PROBLEM 3 %%%%%
\newpage
\begin{problem}[3. (15 points)]
Given that the force of interest is $\delta_t = \frac{2t}{1+2t^2}$,
    \begin{enumerate}[label=(\alph*)]
        \item (9 points) Find the accumulation function $a(t)$.
        \item (6 points) Find the value at time 4 of \$800 to be paid at time $t=9$. Round to 4 decimal places.
    \end{enumerate}
\end{problem}

\newpage

%%%%% PROBLEM 4 %%%%%
\newpage
\begin{problem}[4. (12 points, 4 each)]
Given $d^{(8)} = 12\%$, find the following equivalent rates.


\begin{enumerate}
\itemsep = .01em
    \item[(a)] The effective annual interest rate. 
    \item[(b)] The force of interest $\delta$.
    \item[(c)] The nominal interest rate convertible quarterly.
\end{enumerate} 
\end{problem}

\newpage

%%%%% PROBLEM 5 %%%%%
\newpage
\begin{problem}[5. (16 points)]
Aubie loans Chris \$20,000. Chris repays the loan by paying Aubie \$18,000 at the end of one year and \$2,500 at the end of both three and four years. The money received at $t=1,3$ is immediately reinvested at annual effective interest rate of 3\%.
    \begin{enumerate}[label=(\alph*)]
        \item (8 points) Write out the equation to solve for Chris’s annual effective interest rate. Use the financial calculator to find out the value. Round to 4 decimal places as a percentage.
        \item (8 points) Calculate Aubie's annual yield rate. Round to 4 decimal places as a percentage.
    \end{enumerate}
\end{problem}

\newpage

%%%%% PROBLEM 6 %%%%%
\newpage
\begin{problem}[6. (12 points)]
To pay your tuition you take out a loan of \$25,000. The loan is governed by compound annual discount rate of 5\%. you repay $\$1,000$ at the end of one year, $X$ at the end of three years, \$6,000 at the end of four years and \$10,000 at the end of six years. Calculate $X$. Round to 4 decimal places.
\end{problem}

\newpage

%%%%% PROBLEM 7 %%%%%
\newpage
\begin{problem}[7. (12 points)]
Ellen receives \$2,000 every year on her eighteenth birthday through her thirtieth birthday. using an annual effective discount rate of 5\%. Find the value of her payments on her thirtieth birthday. Round to 2 decimal places. (\textit{Simply entering numbers into a calculator without showing your work will not receive full credit})
\end{problem}

\newpage

%%%%% PROBLEM 8 %%%%%
\newpage
\begin{problem}[8. (12 points)]
Ava makes monthly deposits of $X$ for 5 years with 12\% nominal monthly interest rate. These deposits are exactly enough to allow her to make 10 annual withdrawals of \$1,000. The first withdrawal is one year after the last deposit. The nominal monthly interest rate decreases to 6\% directly after the time of the last deposit and continues forever. Find $X$. Round to 2 decimal places. (\textit{Simply entering numbers into a calculator without showing your work will not receive full credit})
\end{problem}


\end{document}
